% sec-02: the action register and the single approval step.
% Figure 2 (graphviz-type records), Table 2, and the actionbox.
\section{The Action Register}

Five mechanisms hold an inquiry from this company, sent between July 10 and
August 8, 2026 \cite{tenapps2026}. Every one of those inquiries was written
before the approval. This section is the register of what is sent back, to whom,
and what each letter newly asks.

\subsection{Why none of these is a cold letter}

A re-contact after a material change is a normal and welcome event for a program
officer. A fresh introduction five weeks after the first letter reads as though
the first letter was forgotten, and it invites the answer the first letter
already received. So each of the five opens by naming the earlier inquiry and its
date, states the one fact that changed, and asks a question the earlier letter
could not have asked. None re-explains the program from the beginning.

\begin{appfloat}[!tb]
\begin{appfig}[graphviz-type / record nodes in a dashed cluster]
% Five-row, four-column record table.  Column widths are cut to the longest real
% cell in each column.  Row pitch 0.78cm; header row 0.72cm deep.
\foreach \x/\w/\lab in {0/30/{Mechanism}, 2.55/17/{Contacted}, 4.30/25/{Ask}, 6.85/34/{State now}}{
  \node[gvcellh,minimum width=\w mm,text width=\w mm] at (\x,0) {\lab};}
\foreach \y/\m/\d/\a/\s in {%
  -0.78/{NIH SEED, SBIR}/{Aug 2026}/{\$306K Phase I}/{Re-contact sent},
  -1.56/{ARPA-H mission office}/{Aug 2026}/{\$2.1M, three gates}/{Gate 1 narrowing proposed},
  -2.34/{NCI CTEP}/{Aug 2026}/{Concept review}/{Routing question open},
  -3.12/{NIH Pioneer, DP1}/{Jul 2026}/{\$700K per year}/{Codes and effort open},
  -3.90/{NSF TIP, X-Labs}/{Aug 2026}/{\$700K per year}/{Held for day 5}}{
  \node[gvcell,minimum width=30mm,text width=30mm] at (0,\y) {\m};
  \node[gvcell,minimum width=17mm,text width=17mm] at (2.55,\y) {\d};
  \node[gvcell,minimum width=25mm,text width=25mm] at (4.30,\y) {\a};
  \node[gvcells,minimum width=34mm,text width=34mm] at (6.85,\y) {\s};}
\node[gvctitle,anchor=west] (ct) at (-1.62,0.62)
  {Federal mechanisms holding an inquiry from ChemicalQDevice};
% The fit value is braced: its coordinates carry commas, and an unbraced value
% would be split into separate keys by the TikZ key parser.
\node[gvcluster,fit={(ct)(-1.65,0.34)(8.62,-4.28)},inner sep=6pt] {};
\node[gvboxk,text width=44mm] (apr) at (3.50,-5.55)
  {FDA approval, August 26, 2026: one field changes in every record};
\foreach \y in {-0.78,-1.56,-2.34,-3.12,-3.90}{
  \draw[gvedgeb] (apr.north) .. controls (3.50,-4.80) .. (-1.52,\y);}
\end{appfig}
\vspace{-0.60cm}
\figcaption{Figure 2. Five federal mechanisms as records, each with what was sent, when,\\
what was asked, and the state of the file after one change of external fact.}
\end{appfloat}

\subsection{The five letters}

\begin{apptable}
\begin{tabularx}{\textwidth}{>{\raggedright\arraybackslash}p{0.5cm} >{\raggedright\arraybackslash}p{3.4cm} >{\raggedright\arraybackslash}X}
\headrow \thead{\#} & \thead{Recipient} & \thead{The one question this letter newly asks}\\
\midrule
1 & NIH SEED and SBIR program staff & Does an approved agent change the Phase I feasibility bar applied to a combined drug, device, and software workflow \\
2 & ARPA-H mission office & Can gate 1 be narrowed in scope and repriced, since half of what it was written to retire has resolved externally \\
3 & NCI Cancer Therapy Evaluation Program & Is perioperative use of a newly approved agent inside or outside the investigator-initiated concept pathway \\
4 & NIH Common Fund and CSR review contacts & Do the tentative science codes and the 51 percent effort expectation change for a firm rather than an institution \\
5 & The company's broker-dealer & Execution of a four-rung Treasury ladder cut to a nine-month operating horizon \\
\end{tabularx}
\end{apptable}
\vspace{-0.60cm}
\tabcap{Table 2. The five letters of this day, the recipient of each, and the\\
single question in each that the earlier inquiry could not have asked.}

Letter 2 is the one worth pausing on. It proposes a \textbf{reduction} in a filed
outline rather than an expansion of it. Gate 1 of the August outline was written
to retire agent uncertainty alongside workflow uncertainty; half of what it was
for has since been retired externally at no cost to that office. Telling a
program manager this before being asked is cheaper than being asked, and it is
the kind of fact a program manager remembers about an applicant.

Letter 5 belongs in the same register as the other four for a reason that is
easy to miss. It is addressed to a person, it carries a subject line, and it
asks for an action, which is what makes it correspondence rather than a note to
self. It is also the only letter of the five whose recipient can act on the same
day it is sent. The four federal letters open a conversation whose next move
belongs to someone else; the brokerage letter closes one.

\subsection{What is deliberately not sent today}

Two things that could be sent today are not. No approach is made to the agent's
developer, because a developer reasonably asks who is funding the work and the
four federal letters above are the answer to that question; that approach is day
2. No approach is made to a clinical site, because a site asks the same question
in a different form and then asks who the investigator is; that approach is day
3. Sending all three classes of letter in one day would produce three
conversations that each need the other two to have finished first.

\subsection{The single approval step}

\actionbox{One decision is open on this day}{Approve sending the five letters in
\appfile{funding/auto-fund/02Sep26/emails}, and approve the Treasury instruction
in \appfile{funding/auto-fund/02Sep26/investing}. Every address, subject line,
body, attachment manifest, order type and limit is written. No drafting remains.
Letters 1 to 4 and letter 5 should not be sent in the same session: a brokerage
desk and a program office should never appear in the same thread.}
